\chapter{绪论}\label{chap:intro}

\section{研究背景与问题定义}
本文聚焦 PostgreSQL 在 2022--2025 年间的提交演化，关注其中具有代表性的 Bug 修复行为。研究问题是：在真实工程中，Bug 修复是否存在可归纳的结构模式，以及这些模式能否通过静态、动态和形式化三类证据被一致支持。

\section{研究目标与贡献}
本文的目标是构建一个可复现的数据与分析流程，并围绕“识别--分析--验证”形成完整证据链。主要贡献包括：
\begin{itemize}
    \item 构建 PostgreSQL Bug 修复样本集，并建立可执行的标注规范；
    \item 提炼修复结构演化规律，分析控制流与条件复杂度变化；
    \item 在典型案例上结合动态执行轨迹与形式化约束验证修复正确性。
\end{itemize}

\section{研究方法与技术路线}
技术路线分为三层：
\begin{enumerate}
    \item 数据层：提交抓取、规则识别、人工复核与统计分析；
    \item 证据层：AST/pycparser 静态分析与运行轨迹复现；
    \item 验证层：基于 Z3 的状态抽象与可满足性检验。
\end{enumerate}

\section{论文结构}
全文共七章：第二章介绍数据集和识别流程；第三章给出静态结构演化结果；第四章给出动态案例复现；第五章进行形式化验证；第六章综合讨论三类证据及治理启示；第七章总结并展望。
